\documentclass[11pt]{article}
\usepackage[margin=1in]{geometry}
\usepackage{amsmath,amssymb}
\usepackage{xcolor}
\usepackage[hidelinks]{hyperref}
\usepackage{booktabs}

\newcommand{\code}[1]{\texttt{#1}}
\title{\textbf{Hierarchical Semantic GMM}\\[3pt]
\large A coarse code-centered Gaussian mixture for \emph{concepts},\\ with a fine dual branch for \emph{detail}}
\date{}

\begin{document}
\maketitle
\vspace{-2.5em}

\begin{abstract}
\noindent We add a \textbf{coarse, large-receptive-field quantization level} to the DUALVAE whose
codes summarize object-scale regions (concepts) rather than $8{\times}8$ patches (textures), and we
give \emph{that} level the code-centered Gaussian-mixture prior. The existing fine VQ~$+$~continuous
branch is demoted to a \emph{detail/texture} path and encodes the \emph{residual} after the coarse
layout, so the coarse codes are forced to carry the dominant structure while pixel fidelity is
preserved. The result is ``the GMM, moved to the semantic level'': a small, separable, probabilistic
concept codebook on top of a high-fidelity detail reconstructor. This note explains the motivation,
the exact architecture, why it should keep reconstruction, why it is novel, and how to run and
evaluate it.
\end{abstract}

\section{The problem it solves}
Our fine codebook quantizes a $32{\times}32$ latent grid, so each code sees only an $8{\times}8$-pixel
patch. A code with an $8{\times}8$ receptive field \emph{can only} describe local appearance --- so the
codes are a \textbf{texture vocabulary}: a single image shatters into $\sim$$210$ distinct codes,
neighbouring patches share a code only $\sim$$7\%$ of the time, and the bag-of-codes linear probe
plateaus at $\approx0.44$. Cross-view invariance (CVI) and masked code modeling (MCM) tried to make
these codes semantic \emph{without changing the mechanism}, and both barely moved the code map --- a
per-patch codebook has no notion of an object, so no auxiliary loss can turn it into one.

\textbf{The root cause is receptive field.} The fix is therefore structural, not another loss:
model concepts at a scale where a code \emph{can} be a concept.

\section{The idea in one line}
Add a \textbf{coarse level} (downsample the latent again, e.g.\ $\times4$ to an $8{\times}8$ grid, so
each coarse code sees a large region), carry \textbf{our code-centered GMM at that coarse level}
(coarse code $=$ mixture mean, coarse offset $\sim\mathcal{N}(0,\sigma_k^2)$), and let the \textbf{fine
VQ~$+$~continuous branch model only the residual} left after the coarse layout. Coarse $=$ concepts;
fine~$+$~continuous $=$ detail.

\section{Architecture (exact forward)}
Let $z_e^{\text{vq}}\in\mathbb{R}^{B\times C\times 32\times 32}$ be the current pre-quantization latent.
The hierarchical forward (flag \code{hierarchical\_semantic}) is:
\begin{align}
\textbf{coarse encode:}\quad & u = \mathrm{CoarseDown}(z_e^{\text{vq}}) \in \mathbb{R}^{B\times C\times g\times g},\quad g=32/f,\\
\textbf{coarse quantize:}\quad & e_{k}, \;\text{idx},\;\text{commit}_c = \mathrm{VQ_{coarse}}(u),\\
\textbf{coarse GMM offset:}\quad & \Delta_c = \mu_c + \sigma_c\odot\varepsilon,\quad \mu_c,\log\sigma_c^2 = \mathrm{Head}_c(u-e_k^{\text{detach}}),\\
\textbf{coarse latent:}\quad & c = e_k + \Delta_c \sim \textstyle\sum_j \pi_j\,\mathcal{N}(e_j,\sigma_j^2 I)\ \ \text{(the SEMANTIC GMM)},\\
\textbf{broadcast:}\quad & \hat{c} = \mathrm{CoarseUp}(c) \in \mathbb{R}^{B\times C\times 32\times 32},\\
\textbf{fine residual:}\quad & z_e^{\text{fine}} = z_e^{\text{vq}} - \hat{c}^{\,\text{detach}},\\
\textbf{fine path (unchanged):}\quad & z_{\text{vq}},\ \Delta = \text{fine VQ}(z_e^{\text{fine}}) + \text{continuous}(z_e^{\text{fine}}-z_{\text{vq}}^{\text{detach}}),\\
\textbf{decode:}\quad & \hat{x} = \mathrm{Decoder}\big(\mathrm{Attn}(\ \hat{c} + z_{\text{vq}} + \Delta\ )\big).
\end{align}
The coarse KL is the same code-centered form as the fine one, $\mathrm{KL}\!\big(\mathcal{N}(\mu_c,\sigma^2)\,\|\,\mathcal{N}(0,\sigma_{k}^2)\big)$, weighted by \code{coarse\_kl\_beta}; the coarse commitment ties $u$ to its code. Losses:
\[
\mathcal{L} = \underbrace{\mathcal{L}_{\text{recon}} + \mathcal{L}_{\text{GAN}} + \text{fine VQ} + \beta\,\mathrm{KL}_{\text{fine}}}_{\text{unchanged}}
\;+\; \underbrace{\text{commit}_c + \beta_c\,\mathrm{KL}_{\text{coarse}}}_{\text{new}} .
\]

\paragraph{Two design choices that make it work.}
\begin{itemize}
  \item \textbf{Fine encodes the residual $z_e^{\text{vq}}-\hat{c}$.} This is the VQVAE-2 / RQ principle
  applied \emph{across scales}: the coarse level must explain the bulk (or the residual the fine level
  sees is large and reconstruction suffers), which \emph{forces the coarse codes to be structural}.
  The \code{detach} stops the fine level from pushing the coarse codes around.
  \item \textbf{Zero-initialised \code{CoarseUp}.} Its last conv starts at $0$, so $\hat{c}=0$ at step~0:
  the fine level sees the untouched $z_e^{\text{vq}}$ (matching the $k$-means codebook init) and the
  coarse contribution \emph{ramps in} as reconstruction gradient trains \code{CoarseUp}. A clean warm
  start, no early instability.
\end{itemize}

\section{Why reconstruction is preserved}
Semantic-token and object-centric models usually \emph{lose} reconstruction, because a few coarse
tokens cannot hold pixel detail. Here they do not have to: the \textbf{fine VQ~$+$~continuous branch is
the detail channel}, and it is untouched. The coarse level only supplies the low-frequency,
object-scale layout; everything the coarse tokens cannot represent flows into the fine residual and the
continuous offset, exactly as in the flat model (which already reconstructs at rFID~$\approx4$, tying a
matched VAE). So we expect \textbf{concepts at the coarse level at no reconstruction cost}.

\section{Why this is novel}
It is a novel \emph{synthesis} of three ingredients that are individually known but, to our knowledge,
have not been combined:
\begin{enumerate}
  \item \textbf{Code-centered GMM prior at a semantic scale.} Hierarchical VQ (VQVAE-2, RQ-VAE) stacks
  \emph{plain} VQ levels --- discrete tokens with no probabilistic structure. We instead place our
  \emph{code-centered Gaussian-mixture} prior at the coarse level, so the semantic latent is a
  \textbf{separable, probabilistic mixture} (codes as means, per-code variance $\sigma_k^2$), not a bag
  of tokens. That makes the semantic level analysable (separability, code-anchored likelihood) and
  directly \emph{samplable} as a mixture.
  \item \textbf{A dual continuous branch that removes the semantics-vs-fidelity trade-off.} The reason
  hierarchical-VQ / object-centric semantic layers hurt reconstruction is that the semantic bottleneck
  must also carry detail. Our continuous branch is precisely the detail channel those methods lack, so
  the coarse level is \emph{free to be purely semantic}. Getting object-scale concepts \emph{and}
  high-fidelity reconstruction from the same autoencoder is the contribution.
  \item \textbf{A built-in two-level generative prior.} The coarse GMM is a small
  ($g{\times}g$, $K_c$-way) structural prior --- exactly the ``establish a semantic layout, then refine
  detail'' recipe used to condition latent diffusion / flow matching --- but here it is \emph{inside}
  the AE and carries the mixture. A cheap AR or flow prior over the coarse codes plus the analytic
  per-code Gaussians yields a full generative model; the fine~$+$~continuous path renders it.
\end{enumerate}
In short: prior work gives you \emph{either} a hierarchical discrete latent \emph{or} a probabilistic
mixture \emph{or} high-fidelity reconstruction. This design aims for all three at once, by putting the
GMM at the concept scale and delegating pixels to the dual branch.

\section{How to run it}
Gated on \code{hierarchical\_semantic} (default \code{false}), so every existing run is bit-identical.
On the cluster:
\begin{quote}
\code{python main.py --model dualvae --config configs/dualvae\_hier\_semantic\_D.yaml}
\end{quote}
It is Run~D+ (code-centered prior, $\beta{=}0.001$, LPIPS+GAN) plus the coarse block:
\code{coarse\_factor: 4} ($32{\to}8$), \code{coarse\_num\_embeddings: 64} (small $=$ concept-like),
\code{coarse\_kl\_beta: 0.001}. It adds $<\!1$M parameters.

\paragraph{What to watch (wandb).}
\begin{itemize}
  \item \textbf{Train/Coarse Perplexity} --- effective number of concepts in use (out of $64$). Want it
  to settle to a healthy value ($\gtrsim$ a few tens) \emph{without} collapsing to $1$--$2$.
  \item \textbf{Val/rFID} --- must stay $\approx$$4$. If it rises, the coarse level is stealing detail;
  lower \code{coarse\_kl\_beta} or raise \code{coarse\_num\_embeddings}.
  \item \textbf{Train/Coarse KL, Train/Coarse Commitment} --- healthy, non-zero, not exploding.
\end{itemize}

\section{How to evaluate the hypothesis}
The whole point is: \emph{are the coarse codes finally semantic, at no reconstruction cost?} Two direct tests:
\begin{enumerate}
  \item \textbf{Probe the coarse bag-of-codes.} Extract the $64$-d coarse code histogram per image and
  run the frozen linear probe / retrieval. \textbf{Success $=$ beating the fine $0.44$ ceiling.}
  \item \textbf{Re-render the code-cluster overlay on the coarse codes.} If the coarse codes form
  \emph{object regions} (few, large, contiguous) instead of the fine texture speckle, the receptive-field
  hypothesis holds. Compare distinct-codes/image and neighbour-coherence against the fine map.
\end{enumerate}
And confirm \textbf{rFID} is unchanged. If the coarse codes probe higher \emph{and} rFID holds, the
hierarchical semantic GMM did what CVI/MCM could not: it made the codes concepts without paying pixels.

\end{document}
